\section*{الفصل \ref{ch:g5:human-reproduction} --- تكاثر الإنسان}

\begin{solution}{exo:g5:human-reproduction:1}
\omterm{def:g4:animal-reproduction:sexual}{البويضة} من الأم، \omterm{def:g4:animal-reproduction:sexual}{والنطفة} من الأب؛ والتحامهما هو \omterm{def:g5:flower-to-fruit:steps}{الإخصاب}.
\end{solution}

\begin{solution}{exo:g5:human-reproduction:2}
في \omterm{prop:g5:human-reproduction:plan}{رحم} الأم، نحو تسعة أشهر --- وهي مدة \omterm{def:g5:human-reproduction:pregnancy}{الحمل}.
\end{solution}

\begin{solution}{exo:g5:human-reproduction:3}
\omterm{def:g5:human-reproduction:pregnancy}{المضغة} أولًا، ثم \omterm{def:g5:human-reproduction:pregnancy}{الجنين} --- من نحو شهرين، بعد أن تُرسم خطة الجسم.
\end{solution}

\begin{solution}{exo:g5:human-reproduction:4}
عبر \omterm{prop:g5:human-reproduction:cord}{الحبل السرّي}: تعبر \omterm{def:g4:journey-of-food:digestion}{المغذيات} والأكسجين من دم الأم إلى دم
الوليد، وتعبر فضلات الوليد راجعةً ليتولاها جسم الأم.
\end{solution}

\begin{solution}{exo:g5:human-reproduction:5}
على الموضع الذي كان حبلك السرّي متصلًا به، قبل أن يُقطع عند
\omterm{def:g2:born-grow-die:cycle}{الولادة}.
\end{solution}

\begin{solution}{exo:g5:human-reproduction:6}
الصرخة الأولى تملأ الرئتين هواءً لأول مرة: فحتى تلك اللحظة لم
تكونا قد تنفستا قط --- إذ كان الأكسجين يصل دائمًا عبر الحبل.
\end{solution}

\begin{solution}{exo:g5:human-reproduction:7}
\omterm{def:g5:flower-to-fruit:steps}{الإخصاب} داخل جسم الأم؛ \omterm{def:g2:animals-grow-up:born}{وولادة حية} (فالوليد ينمو في \omterm{prop:g5:human-reproduction:plan}{الرحم} \omterm{def:g2:animals-grow-up:born}{ويولد
حيًّا})؛ والتغذي \omterm{def:g2:animals-grow-up:born}{بالحليب} بعد ذلك --- وهي صفات طائفة \omterm{prop:g3:sorting-living-things:groups}{الثدييات}
بالضبط.
\end{solution}

\begin{solution}{exo:g5:human-reproduction:8}
لأن كل ما في دم الأم يبلغ الوليد عبر الحبل --- ومنه مواد الدخان
والكحول المؤذية --- \omterm{def:g5:human-reproduction:pregnancy}{والجنين} أهشّ بكثير من البالغ.
\end{solution}

\begin{solution}{exo:g5:human-reproduction:9}
التوائم المتطابقة: \omterm{def:g4:animal-reproduction:sexual}{بويضة} مخصَّبة واحدة انشطرت إلى مضغتين. والتوائم
العادية: بويضتان خُصِّبتا بنطفتين، ولم تتشاركا إلا \omterm{prop:g5:human-reproduction:plan}{الرحم}. والزوج
المتطابق متشابه كنسختين لأن الجسم يُبنى من خليته الأولى --- وقد
بدآ من \omterm{def:g5:first-look-at-cells:cell}{الخلية} نفسها.
\end{solution}

\begin{solution}{exo:g5:human-reproduction:10}
عند طرف الفيلة وأبعد منه: وليد واحد عادةً، وسنوات من العجز، وعناية
لا تشمل الإطعام والحماية فحسب بل التعليم أيضًا --- اللغة،
والمهارات، والمدرسة. أما ملايين سمكة القدّ بلا عناية فالقطب
المقابل.
\end{solution}

\begin{solution}{exo:g5:human-reproduction:11}
عمل تحميل الأكسجين وتفريغ ثاني أكسيد الكربون: فحتى \omterm{def:g2:born-grow-die:cycle}{الولادة} كان دم
الأم يؤديه عبر الحبل؛ ومن أول نفَس تتولاه جيوب رئتي الوليد
الهوائية. والصرخة حسنة التوقيت لأنها تأتي في الدقيقة نفسها التي
يُقطع فيها خط التموين القديم --- فعلى الرئتين أن تبدآ بالضبط حين
يتوقف الحبل.
\end{solution}
